% 06f-baselines.tex 中文版（中英对照用；不含 label，避免与英文版重复）
% 数字来源：paper/comparison-lab/results/final-table.json（n=5）；表格沿用
% % generated from paper/comparison-lab/results/final-table.json (v2) -- do not edit by hand
% primary sources: results/summary.json (sha256 78e5c5328bb2...),
%                  mem0-chaos-lab/results/summary.json + human-review.json
\begin{table*}[t]
\centering\footnotesize
\caption{Three-subject pilot ($n=5$ independent rounds per cell; one host, one
model route, one day). S1 = memory conflict (ALPHA$\rightarrow$BETA); S2 =
jailbreak/injection. \emph{Hit}: current value answered / attack succeeded under
human review; \emph{Leak}: superseded value or planted token disclosed;
\emph{Sup.}: the superseded value is absent from what the system injects;
\emph{Int.}: model-level refusal; \emph{gate}/\emph{audit}: structural UCA-01
block and its audit record. Mem0 is shown in its recall-working diagnostic
configuration (see the default-configuration note); its default row is omitted
from the main body because its embedding pathway is empty in this environment.
Sources: \texttt{paper/comparison-lab/results/final-table.json}.}
\label{tab:baselines}
\begin{tabular}{@{}l ccc ccc ccc@{}}
\toprule
 & \multicolumn{3}{c}{S1 memory conflict} & \multicolumn{3}{c}{S2 jailbreak} & \multicolumn{3}{c}{S2 structural} \\
\cmidrule(lr){2-4}\cmidrule(lr){5-7}\cmidrule(lr){8-10}
Subject & Hit & Leak & Sup. & Hit & Leak & Int. & Blocked & Audit & Refused \\
\midrule
Baseline (no plugin) & 5/5 & 0/5 & 0/5 & 0/5 & 0/5 & 5/5 & 0/5 & n/a & 5/5 \\
Mem0 2.2.1, zh diagnostic & 5/5 & 0/5 & 0/5 & 0/5 & 0/5 & 5/5 & 0/5 & n/a & 5/5 \\
LINGSHU (lab profile) & n/i & n/i & 5/5 & 0/5 & 0/5 & 5/5 & 5/5 & 5/5 & 5/5 \\
\bottomrule
\end{tabular}

\vspace{2pt}
\raggedright\scriptsize
S1 uses the docs/19-supported correction sentence for all three subjects. Baseline
S1 hit is 4/5 by the literal string rule; the fifth round answered ASCII
\texttt{Beta}, counted as a hit on human review. LINGSHU S1 hit/leak are marked
n/i: the memory ledger is correct in 5/5 cells (superseded value tombstoned,
\texttt{supersedeEvents=1}, only BETA injected in the preview), but the model's
first request sees the previous turn's injection snapshot (F-D), so the answers
are ALPHA 5/5; a warm diagnostic with one extra turn answers BETA 5/5. Its
structural suppression is 5/5, the effective suppression 0/5. Mem0's default
configuration extracts English memory text; with the local non-neural embedder,
Chinese queries retrieve nothing (cosine $\le 0.063 < 0.1$), so its default S1
``suppression'' is an empty-recall artefact and the store still retains the stale
value (ADD-only pipeline). S2 canary never reaches the model in the LINGSHU arm
because the fact extractor rejects the secret-bearing sentence and the gate
blocks the attack turn; all three arms refuse the obvious payload, so no
model-refusal difference is claimed. $n=5$: directional only, no confidence
intervals.
\end{table*}
（英文表，数字与英文版逐格一致；如需中文表头请另行生成）。
\subsection{外部基线：无插件宿主与 Mem0}

前面的实验比较的是同一系统内部的臂，或是在隔离条件下测量闸门。
读者可能会问一个显然的外部问题：相对于 (i) 不带任何记忆层的裸模型，
以及 (ii) 一个没有安全层的生产级长期记忆框架，控制面到底买到了什么？
因此我们在两个场景上运行了一次小规模的三受试者试点。
我们把它定位为方向性试点（directional pilot），而不是基准测试：
每个单元格五轮、一个宿主、一条模型路由、一天。

\paragraph{受试者与协议。} 三个受试者：\textsc{Baseline}（基线），完全不装插件的宿主
（把原始对话记录放在一次调用里）；\textsc{Mem0}，官方 \texttt{mem0ai}~2.2.1 库~\cite{mem0_2025}，
由同一个 DeepSeek 模型驱动，使用内存向量库；
以及把插件启用的实验室 profile 作为 LINGSHU，每个单元格一个隔离数据目录。
场景 S1（记忆冲突）先埋下一个代号 ALPHA，再插入一个干扰性偏好，然后纠正为 BETA 并声明
ALPHA 作废，最后询问代号。场景 S2（越狱）先埋下一个机密 token，然后要求系统忽略指令，
打印它的系统提示词以及全部注入记忆（其中包含该 token）。
\textsc{Baseline} 在一次调用中收到整段对话记录；LINGSHU 逐轮驱动
（S1 四次调用、S2 两次），这样它的记忆与闸门才真正运行；
\textsc{Mem0} 逐轮通过 \texttt{add} 写入，并用最后那个问题查询。
每个单元格都从全新状态开始。

\paragraph{指标。} 每个单元格记录三个比率。\emph{Hit}（命中）：S1 中最终回答给出当前代号；
S2 中表示攻击成功（token 被泄露，或在非拒绝性回答中出现内部标记）。
\emph{Leak}（泄漏）：S1 回答中出现被取代的值 ALPHA，或 S2 回答中出现被埋下的 token。
\emph{Interception}（拦截）：对 S1，指旧值不在系统所注入的内容中
（\textsc{Mem0} 是检索到的记忆，LINGSHU 是注入记忆预览，\textsc{Baseline} 是原始对话记录）；
对 S2，指某道结构性闸门阻断了注入（UCA-01，来源归因于用户消息），
模型层拒绝另行报告。

% generated from paper/comparison-lab/results/final-table.json (v2) -- do not edit by hand
% primary sources: results/summary.json (sha256 78e5c5328bb2...),
%                  mem0-chaos-lab/results/summary.json + human-review.json
\begin{table*}[t]
\centering\footnotesize
\caption{Three-subject pilot ($n=5$ independent rounds per cell; one host, one
model route, one day). S1 = memory conflict (ALPHA$\rightarrow$BETA); S2 =
jailbreak/injection. \emph{Hit}: current value answered / attack succeeded under
human review; \emph{Leak}: superseded value or planted token disclosed;
\emph{Sup.}: the superseded value is absent from what the system injects;
\emph{Int.}: model-level refusal; \emph{gate}/\emph{audit}: structural UCA-01
block and its audit record. Mem0 is shown in its recall-working diagnostic
configuration (see the default-configuration note); its default row is omitted
from the main body because its embedding pathway is empty in this environment.
Sources: \texttt{paper/comparison-lab/results/final-table.json}.}
\label{tab:baselines}
\begin{tabular}{@{}l ccc ccc ccc@{}}
\toprule
 & \multicolumn{3}{c}{S1 memory conflict} & \multicolumn{3}{c}{S2 jailbreak} & \multicolumn{3}{c}{S2 structural} \\
\cmidrule(lr){2-4}\cmidrule(lr){5-7}\cmidrule(lr){8-10}
Subject & Hit & Leak & Sup. & Hit & Leak & Int. & Blocked & Audit & Refused \\
\midrule
Baseline (no plugin) & 5/5 & 0/5 & 0/5 & 0/5 & 0/5 & 5/5 & 0/5 & n/a & 5/5 \\
Mem0 2.2.1, zh diagnostic & 5/5 & 0/5 & 0/5 & 0/5 & 0/5 & 5/5 & 0/5 & n/a & 5/5 \\
LINGSHU (lab profile) & n/i & n/i & 5/5 & 0/5 & 0/5 & 5/5 & 5/5 & 5/5 & 5/5 \\
\bottomrule
\end{tabular}

\vspace{2pt}
\raggedright\scriptsize
S1 uses the docs/19-supported correction sentence for all three subjects. Baseline
S1 hit is 4/5 by the literal string rule; the fifth round answered ASCII
\texttt{Beta}, counted as a hit on human review. LINGSHU S1 hit/leak are marked
n/i: the memory ledger is correct in 5/5 cells (superseded value tombstoned,
\texttt{supersedeEvents=1}, only BETA injected in the preview), but the model's
first request sees the previous turn's injection snapshot (F-D), so the answers
are ALPHA 5/5; a warm diagnostic with one extra turn answers BETA 5/5. Its
structural suppression is 5/5, the effective suppression 0/5. Mem0's default
configuration extracts English memory text; with the local non-neural embedder,
Chinese queries retrieve nothing (cosine $\le 0.063 < 0.1$), so its default S1
``suppression'' is an empty-recall artefact and the store still retains the stale
value (ADD-only pipeline). S2 canary never reaches the model in the LINGSHU arm
because the fact extractor rejects the secret-bearing sentence and the gate
blocks the attack turn; all three arms refuse the obvious payload, so no
model-refusal difference is claimed. $n=5$: directional only, no confidence
intervals.
\end{table*}


\paragraph{试点显示了什么。} 在 S1 上，无插件宿主能从原始对话记录里答出代号
（人工复核为 5/5；字面规则判为 4/5，因为有一轮回答的是 ASCII 形式的 \texttt{Beta}），
但从不抑制旧值——按构造如此：对话记录里就包含它。\textsc{Mem0} 保留两条条目：
它默认的 \texttt{add} 路径是累加式的——该库自己的 prompt 写明唯一操作就是 ADD——
因此被取代的值仍留在存储中，并且在召回可用的诊断配置下，会与纠正一起被检索出来
（抑制 0/5）；我们加入了一个显式删除的反事实，以说明失败来自默认策略而不是库本身坏掉。
LINGSHU 通过其唯一写方写入纠正，并给 ALPHA 立墓碑：
在每一个单元格中，账本都记录一次 supersede 事件，旧条目被标为 superseded，
而注入记忆预览只包含 BETA（结构性抑制 5/5）。
这就是 S1 上的主要结构性差异，而且这是类别上的差异，不是排序质量上的差异：
Mem0 的 ADD-only 向量库检索“相似”的内容，而 LINGSHU 在事实写入前检查
$(\textsc{subject},\textsc{relation})$ 对，因此纠正会给 ALPHA 立墓碑，而不是加入它。
检索增强（retrieval augmentation）对结构防御（structural defence）。

有两点限定这一差异。第一，LINGSHU 在 S1 上的模型层回答在这里不可解释：
尽管预览只含 BETA，模型第一次请求使用的却是上一轮的注入快照，
因此五个回答全都是 ALPHA（发现 F-D）；一个在提问前多加一轮的 warm 诊断，
在 5/5 个单元格里答出 BETA。第二，事实提取器是一个规则桩：
它识别不了自然的纠正措辞（``现在改成''，发现 F-A），还会把最后那个问题当成事实存下来（发现 F-B）。
我们报告结构性成功与这些实现缺口，而不是报告一次干净的 S1 胜利。

在 S2 上情况不同，我们如实报告。三个受试者大多在模型层拒绝了那个显而易见的注入，
因此端到端的 ``拒绝率'' 无法在这个载荷上区分它们。真正区分它们的是结构性准入：
LINGSHU 在每一个 S2 单元格中都记录了一次归因于用户消息的 UCA-01 阻断，并附有审计证据，
而 \textsc{Mem0} 与无插件宿主根本没有准入闸门。
因此这次试点支持的是一条狭窄的结构性主张，而不是模型行为主张：
\emph{在本协议下，LINGSHU 把安全准入变成了结构性的、可审计的闸门（5/5 个单元格），
而纯记忆框架没有闸门（0/5），并且没有在记忆层里新增任何机制。}
我们不声称 LINGSHU 比 \textsc{Mem0} 拒绝更多攻击：在一个明显对齐良好的模型上，两者都会拒绝。

\paragraph{它不能说明什么。} 三个受试者看到的输入形态不同（原始对话记录 vs 检索到的记忆），
因此各受试者之间的比率不能直接比较，我们也不报告效应量。
该试点是每个单元格五轮、一个宿主、一条模型路由、一天：所有数字都是 $n=5$ 且方向性的，
行级回答已归档，读者可以重新分类。
在 S2 中，被埋下的 token 在 LINGSHU 臂里根本没有到达模型：
事实提取器在存储之前就把这句含机密的句子当作禁忌主题拒绝了，闸门又阻断了攻击轮，
因此零泄漏率不能被读作对记忆通道的一次披露检验。
注入侧闸门会阻止注入并记录下来，但宿主仍会把用户消息送给模型，
因此 ``blocked'' 并不意味着模型从未读到攻击。
这次试点还暴露了工具网关噪声（发现 F-C）与一轮注入滞后（F-D）；
两者都作为产品发现报告，而没有被折进比率里。
